\section{研究议程与实验蓝图}

本章把课堂观点进一步转成可立项的研究议程。目标不是追求概念完整性，而是给出可在 4--12 周内复现实验信号的方案。

\subsection{议程一：协作分布与对抗分布的联合优化}
研究问题：同一模型同时服务正常用户和攻击用户时，如何避免一端优化拉垮另一端？

\begin{knowledgebox}{建议实验设置}
\begin{itemize}
    \item 构建双数据池：$D_c$（collaborative episodes）与 $D_a$（adversarial episodes）。
    \item 训练时按比例混采：$D = \lambda D_c + (1-\lambda)D_a$，并周期性扫描 $\lambda$。
    \item 指标分开报：协作成功率 $S_c$、对抗防御率 $R_a$、综合成本 $C$。
\end{itemize}
\end{knowledgebox}

这种设置的重要性在于，它把``安全性''从发布门槛变成训练一等公民。对于课程团队而言，这也能避免仅凭单一 leaderboard 做错误结论。

\subsection{议程二：基于 exploiter 的自动发现与修复循环}
研究问题：如何让系统自动发现新失败模式，并在可控范围内迭代修复？

\begin{importantbox}{最小可运行修复循环}
\begin{enumerate}
    \item 由 exploiter 生成挑战任务并执行攻击。
    \item 将失败轨迹归因到策略错误、工具错误或规则错误。
    \item 按归因类型进行 targeted finetuning 或 policy patch。
    \item 用冻结评测集验证是否真实修复、是否引入新回归。
\end{enumerate}
\end{importantbox}

可用如下伪代码描述循环逻辑：
\begin{lstlisting}[caption={Auto-Exploit to Auto-Repair 循环（示意）}]
for round in training_rounds:
    attack_batch = exploiter.generate(main_agent, task_pool)
    failures = evaluate(main_agent, attack_batch)
    labeled = root_cause_label(failures)  # policy / tool / rule
    main_agent = targeted_update(main_agent, labeled)
    report = regression_eval(main_agent, fixed_benchmark)
    if report.safety_drop:
        rollback(main_agent)
\end{lstlisting}

\begin{warningbox}{议程二的风险点}
如果只关注攻击成功率下降，可能出现``过拟合现有攻击脚本''。必须持续引入新攻击策略，并保留不可见测试集，否则很快形成虚假安全感。
\end{warningbox}

\subsection{议程三：match-making 与学习信号密度}
研究问题：如何在任务极难与极易之间保持稳定学习？

\begin{table}[H]
\centering
\caption{三种 match-making 策略的课堂外推}
\begin{tabular}{p{3.0cm}p{4.0cm}p{5.6cm}}
\toprule
策略 & 预期优点 & 潜在缺点 \\
\midrule
Hardest-first & 快速暴露上限缺陷 & 易出现无梯度区，训练不稳定 \\
50/50-first & 学习信号密度高，收敛平稳 & 暴露极端风险速度较慢 \\
Curriculum hybrid & 兼顾稳定与极端覆盖 & 调度规则复杂，参数敏感 \\
\bottomrule
\end{tabular}
\end{table}

\begin{knowledgebox}{建议的课程实验路线}
先用 50/50-first 建立稳定基线，再逐步注入 hardest-first 任务；最后用 curriculum hybrid 做全局平衡。每步都固定随机种子与任务池版本，确保结果可复现。
\end{knowledgebox}

\subsection{议程四：可控性（controllability）与可解释性}
讲者在 53 分钟附近提到 ``having controllability mattered''。这点在 LLM Agent 中尤其重要，因为任务语言开放、动作后果复杂，可控性直接关联到安全合规。

\begin{importantbox}{可控性设计清单}
\begin{itemize}
    \item 行为前置约束：权限级别、调用白名单、预算上限。
    \item 行为中可观测：每步动作理由、输入输出摘要、风险标签。
    \item 行为后可追溯：轨迹存档、异常聚类、快速回滚。
\end{itemize}
\end{importantbox}

\begin{warningbox}{可解释性不是附加报告}
如果解释数据不进入训练回路，它就只是审计文本。真正有价值的解释，必须能反向驱动数据再采样、奖励修正和策略更新。
\end{warningbox}

\subsection{本章小结}
本章给出的研究议程可以概括为四件事：联合优化双分布、构建自动修复循环、设计有效匹配策略、把可控性写进训练系统。做到这四点，课堂理念就能从``观点''变成``生产力''。

